\chapter{Results and Evaluation}
\label{ch:results}

\section{Correctness and Stage-Level Summary}

All 12 benchmark cases passed their element-by-element correctness checks
across every optimisation stage. This is the prerequisite for
interpreting any performance figure: a speedup number is only meaningful
if the accelerator produces a correct answer, and the hardware-in-the-loop
correctness check described in Chapter~\ref{ch:methodology} gives
confidence that synthesis and place-and-route have not introduced
functional discrepancies relative to the RTL simulation. The correctness
comparison itself is made against the INT8 CPU reference on the final
mainline build, with the full-precision CPU result computed separately
as an untimed sanity check; accelerator outputs matched their respective
CPU INT8 references exactly for all 12 cases at all three stages.

Table~\ref{tab:stage-summary} summarises the performance of the three
recorded stages at a glance. Two observations stand out immediately.
First, every stage improved over the previous one without any correctness
regression, which is itself a non-trivial outcome in a hardware project
where incremental changes to the FSM and datapath could easily have
introduced off-by-one or sign-extension bugs. Second, the largest gain
came from reducing memory traffic (INT8 packing: $2.18\times$) rather
than from increasing arithmetic parallelism (DSP MAC: $1.32\times$). This
confirms the memory-bandwidth bottleneck hypothesis developed in
Chapter~\ref{ch:motivation} and Chapter~\ref{ch:design}, and it is the
central experimental result of the project.

\begin{table}[htbp]
  \centering
  \small
  \begin{tabular}{lcccc}
  \toprule
  Stage & Avg speedup & Best case & Total accel cycles & Reduction vs prev. \\
  \midrule
  Baseline accelerator & $16.12\times$ & T4: $18.12\times$ & 427,299 & --- \\
  Parallel DSP MAC     & $20.87\times$ & T4: $24.03\times$ & 324,568 & $1.32\times$ \\
  INT8 packed          & $45.28\times$ & T4: $56.80\times$ & 148,737 & $2.18\times$ \\
  \bottomrule
  \end{tabular}
  \caption{Summary of the three staged RISC-V accelerator evaluations
           across the full 12-case benchmark suite.}
  \label{tab:stage-summary}
\end{table}

\section{Baseline Accelerator}

Table~\ref{tab:baseline-results} presents the full baseline accelerator
results against the PicoRV32 software reference. Even the unoptimised
baseline achieves a consistent and substantial speedup, rising from
$12.12\times$ at the smallest case (T1) to $18.12\times$ at T4, which
demonstrates that fixed hardware setup overhead is being amortised over
a growing body of useful computation. T1 is the worst case in the suite
because its $8 \times 8$ matrices are so small that the FSM overhead (the
$M_{\mathrm{max}} = 64$ cycles required to clear the C-tile, the row
loads, and the final writeback) represents a much larger fraction of the
total cycle count than it does for the larger problem sizes. The trend
across T1--T4 is therefore not only a raw speedup curve but also direct
evidence that the structural design choices discussed in
Section~\ref{sec:fsm} carry the expected amortisation behaviour.

\begin{table}[htbp]
  \centering
  \small
  \begin{tabular}{lp{2.6cm}rrrr}
  \toprule
  Test ID & $M \times K \times N$ / Sparsity & CPU cycles & Accel cycles & Speedup & Result \\
  \midrule
  T1  & $8 \times 8 \times 8$, 75\%    & 10,373   &    856 & $12.12\times$ & PASS \\
  T2  & $16 \times 16 \times 8$, 75\%  & 36,573   &  2,318 & $15.78\times$ & PASS \\
  T3  & $32 \times 32 \times 8$, 75\%  & 136,877  &  7,809 & $17.53\times$ & PASS \\
  T4  & $64 \times 64 \times 8$, 75\%  & 529,101  & 29,195 & $18.12\times$ & PASS \\
  T5  & $64 \times 64 \times 16$, 75\% & 986,829  & 58,112 & $16.98\times$ & PASS \\
  T6  & $64 \times 64 \times 32$, 75\% & 1,902,285 & 115,980 & $16.40\times$ & PASS \\
  T7  & $32 \times 32 \times 32$, 50\% & 951,213  & 58,112 & $16.37\times$ & PASS \\
  T8  & $32 \times 32 \times 32$, 75\% & 491,693  & 30,470 & $16.14\times$ & PASS \\
  T9  & $32 \times 32 \times 32$, 90\% & 215,263  & 13,844 & $15.55\times$ & PASS \\
  T10 & $32 \times 32 \times 32$, 75\%, row-skewed & 491,693 & 30,470 & $16.14\times$ & PASS \\
  T11 & $32 \times 32 \times 32$, 75\%, clustered  & 491,693 & 30,470 & $16.14\times$ & PASS \\
  T12 & $64 \times 64 \times 32$, 90\% & 800,155  & 49,663 & $16.11\times$ & PASS \\
  \bottomrule
  \end{tabular}
  \caption{Baseline PicoRV32 software versus baseline accelerator results.
           Speedup ranges from $12.12\times$ to $18.12\times$ across the
           12-case suite.}
  \label{tab:baseline-results}
\end{table}

\section{Parallel DSP MAC Enhancement}

Table~\ref{tab:dsp-results} shows the results after the DSP-oriented MAC
restructuring described in Chapter~\ref{ch:design}. The optimisation
reduces total accelerator cycles from 427,299 to 324,568, a $1.32\times$
improvement, and every individual case records a lower absolute cycle
count than its baseline counterpart. The improvement is consistent but
modest, and this is itself informative: it means that even with
optimally scheduled DSP arithmetic, the accelerator is not
compute-bound. The memory transaction overhead of loading the
$B$-segment ($TN = 8$ reads per nonzero at 32-bit precision) remains
the dominant cost, and no restructuring of the MAC arithmetic alone can
address that.

\begin{table}[htbp]
  \centering
  \small
  \begin{tabular}{lrrrr}
  \toprule
  Test ID & CPU cycles & Accel cycles & Speedup & Result \\
  \midrule
  T1  & 10,373   &    737 & $14.07\times$ & PASS \\
  T2  & 36,573   &  1,859 & $19.67\times$ & PASS \\
  T3  & 136,877  &  6,024 & $22.72\times$ & PASS \\
  T4  & 529,101  & 22,021 & $24.03\times$ & PASS \\
  T5  & 986,829  & 43,781 & $22.54\times$ & PASS \\
  T6  & 1,902,285 & 87,301 & $21.79\times$ & PASS \\
  T7  & 951,213  & 43,781 & $21.73\times$ & PASS \\
  T8  & 491,693  & 23,296 & $21.11\times$ & PASS \\
  T9  & 215,263  & 10,988 & $19.59\times$ & PASS \\
  T10 & 491,693  & 23,296 & $21.11\times$ & PASS \\
  T11 & 491,693  & 23,296 & $21.11\times$ & PASS \\
  T12 & 800,155  & 38,188 & $20.95\times$ & PASS \\
  \bottomrule
  \end{tabular}
  \caption{Results after enabling the parallel DSP-oriented MAC
           enhancement. Total accelerator cycles fell from 427,299 to
           324,568, a $1.32\times$ reduction.}
  \label{tab:dsp-results}
\end{table}

\section{INT8 Packing}

Table~\ref{tab:int8-results} presents the final mainline INT8 results,
which produced by some distance the largest performance step in the
project. Total accelerator cycles fell from 324,568 to 148,737, a
$2.18\times$ reduction, and the mechanism is precisely what the roofline
analysis of Chapter~\ref{ch:motivation} predicted: cutting the
$B$-segment reads per nonzero from 8 to 2 reduces the dominant
memory-bandwidth cost by a factor of four on the reads themselves, and
the overall cycle-count impact is $2.18\times$ because not every cycle
is memory-bound. Row loads, nonzero index fetches, and the final C-tile
writeback contribute fixed overheads that do not shrink with INT8 and
that now account for a larger share of the total, which is the same
amortisation trend seen in the other direction from T1 to T4 in the
baseline.

\begin{table}[htbp]
  \centering
  \small
  \begin{tabular}{lrrrr}
  \toprule
  Test ID & CPU cycles (INT8) & Accel cycles & Speedup & Result \\
  \midrule
  T1  & 10,735   &    567 & $18.93\times$ & PASS \\
  T2  & 38,087   &  1,111 & $34.28\times$ & PASS \\
  T3  & 142,999  &  2,964 & $48.25\times$ & PASS \\
  T4  & 553,655  &  9,747 & $\mathbf{56.80\times}$ & PASS \\
  T5  & 1,044,151 & 19,216 & $54.34\times$ & PASS \\
  T6  & 2,025,143 & 38,171 & $53.05\times$ & PASS \\
  T7  & 1,012,631 & 19,216 & $52.70\times$ & PASS \\
  T8  & 522,391  & 11,039 & $47.32\times$ & PASS \\
  T9  & 227,481  &  6,109 & $37.24\times$ & PASS \\
  T10 & 522,391  & 11,039 & $47.32\times$ & PASS \\
  T11 & 522,391  & 11,039 & $47.32\times$ & PASS \\
  T12 & 849,333  & 18,519 & $45.86\times$ & PASS \\
  \bottomrule
  \end{tabular}
  \caption{Final INT8-packed implementation results. Peak speedup of
           $56.80\times$ on T4; average speedup of $45.28\times$ across
           the 12-case suite. Total accelerator cycles fell from 324,568
           to 148,737, a $2.18\times$ reduction over the DSP stage.}
  \label{tab:int8-results}
\end{table}

\section{Comparison Against ARM Cortex-M0}

Table~\ref{tab:cortex-results} compares the INT8 accelerator cycle counts
against a software SpMM reference running on an ARM Cortex-M0. The
Cortex-M0 cycle counts are derived from the same benchmark definitions
and the known CPI characteristics of the Cortex-M0 ISA (predominantly
single-cycle instructions with 2--3 cycle load latencies for word
accesses), and they represent a substantially stronger software baseline
than PicoRV32 because the Cortex-M0 has a more efficient instruction
pipeline and faster load handling. Speedups are therefore correspondingly
smaller, but the accelerator still achieves up to $11.92\times$ speedup
on T4, a meaningful result for a small FPGA-attached accelerator, and
arguably a more realistic estimate of the practical benefit of the
hardware than the PicoRV32 figure because the Cortex-M0 is closer to the
kind of processor that such an accelerator would plausibly pair with in
a commercial design.

\begin{table}[htbp]
  \centering
  \small
  \begin{tabular}{lrrrr}
  \toprule
  Test ID & Cortex-M0 cycles & Accel cycles & Speedup & Result \\
  \midrule
  T1  & 2,952   &    567 & $5.21\times$  & PASS \\
  T2  & 9,448   &  1,111 & $8.50\times$  & PASS \\
  T3  & 31,912  &  2,964 & $10.77\times$ & PASS \\
  T4  & 116,200 &  9,747 & $\mathbf{11.92\times}$ & PASS \\
  T5  & 207,592 & 19,216 & $10.80\times$ & PASS \\
  T6  & 390,376 & 38,171 & $10.23\times$ & PASS \\
  T7  & 195,240 & 19,216 & $10.16\times$ & PASS \\
  T8  & 104,488 & 11,039 & $9.47\times$  & PASS \\
  T9  & 46,892  &  6,109 & $7.68\times$  & PASS \\
  T10 & 99,748  & 11,039 & $9.04\times$  & PASS \\
  T11 & 102,209 & 11,039 & $9.26\times$  & PASS \\
  T12 & 169,710 & 18,519 & $9.16\times$  & PASS \\
  \bottomrule
  \end{tabular}
  \caption{ARM Cortex-M0 software SpMM versus the final INT8 accelerator.
           The Cortex-M0 is a stronger baseline than PicoRV32, so
           speedups are correspondingly lower; the accelerator still
           achieves up to $11.92\times$ speedup.}
  \label{tab:cortex-results}
\end{table}

\section{Trend Analysis}

\subsection{Scaling with Problem Size (T1--T4)}

Figure~\ref{fig:speedup-size} [placeholder] shows speedup versus matrix
size for the three stages. Speedup rises consistently from T1 to T4 in
all three, and in the INT8 stage it grows from $18.93\times$ at T1
($8 \times 8 \times 8$) to $56.80\times$ at T4 ($64 \times 64 \times 8$).
The increase is sublinear in the problem size: T1 to T2 corresponds to
an eight-fold increase in arithmetic but only a $1.81\times$ speedup
increase, and T3 to T4 corresponds to a four-fold increase in arithmetic
but only a $1.18\times$ speedup increase, which indicates that the
accelerator is beginning to approach the plateau where fixed overhead
becomes a small fraction of the total cycle count. The primary driver of
this behaviour is the ratio of useful MAC work to FSM overhead: for T1
the $8 \times 8$ matrix at 75\% sparsity contains only about 12 nonzeros
in total, and the C-tile clear alone (64 cycles for $M_{\max} = 64$)
dominates the cycle budget, whereas for T4 the roughly 1024 nonzeros
distribute the same fixed setup cost over a much larger body of work.

\begin{figure}[htbp]
  \centering
  \fbox{\parbox[c][55mm][c]{0.88\linewidth}{\centering
  Placeholder: speedup vs matrix size (T1--T4) for all three stages.
  X-axis: matrix dimension ($8, 16, 32, 64$). Y-axis: speedup.
  Three lines: Baseline, DSP MAC, INT8.}}
  \caption{Speedup versus matrix size for tests T1--T4. All three stages
           show increasing speedup as problem size grows, confirming that
           fixed setup overhead is amortised effectively at larger
           problem sizes.}
  \label{fig:speedup-size}
\end{figure}

\subsection{Effect of Dense Width (T4--T6)}

The sequence T4 ($N = 8$), T5 ($N = 16$), T6 ($N = 32$) tests how
speedup changes as more tile columns must be processed at fixed $M = K$.
In the INT8 stage the speedup moves from $56.80\times$ at T4 through
$54.34\times$ at T5 to $53.05\times$ at T6, staying high and falling only
modestly as $N$ grows. The accelerator handles multiple tile columns by
repeating the same row traversal once per tile column, with a tile-clear
and tile-writeback phase separating each, and the CPU software also
scales linearly with $N$ in its inner loop; the near-constant speedup
across T4--T6 therefore reflects the fact that both sides scale
approximately linearly with $N$ and the accelerator maintains its
structural advantage throughout. The slight decrease from $56.80\times$
to $53.05\times$ reflects the increasing relative cost of the
tile-management phases themselves: with four tile columns instead of
one, four separate tile-clear and writeback phases run sequentially, and
the fixed per-tile-column overhead accumulates against the accelerator's
cycle count rather than the CPU's.

\subsection{Effect of Sparsity (T7--T9)}

The sequence T7 (50\%), T8 (75\%), T9 (90\%) fixes $M = K = N = 32$ and
sweeps the sparsity level, which in the INT8 stage gives
$52.70\times \to 47.32\times \to 37.24\times$. Speedup falls as the
matrix becomes sparser, and the reason is structural: at very high
sparsity both CPU and accelerator perform less arithmetic, but the
accelerator's per-row overhead (tile clear, row-bound load) is paid
regardless of whether the row is populated. At 90\% sparsity many rows
of $A$ are empty, the accelerator still visits each of the 32 rows and
contributes fixed traversal cycles, and the CPU simply skips its inner
loop for those empty rows. Even so, the $37.24\times$ achieved at T9 is
a substantial speedup and shows that the design is practical well into
the sparsity regimes typical of GNN workloads. Real social-graph GNN
matrices often exceed 99\% sparsity, where the absolute cycle count is
low and the speedup may continue to fall, but the acceleration in
absolute time would remain valuable.

\subsection{Pattern Sensitivity (T8, T10, T11)}

T8, T10, and T11 share identical dimensions, sparsity, and NNZ, and
differ only in how the nonzeros are distributed: uniform, row-skewed, or
clustered. In the INT8 stage all three report identical accelerator
cycle counts of 11,039 and identical CPU counts of 522,391. This
insensitivity is not accidental. The accelerator processes each row
sequentially regardless of how many nonzeros that row contains, and the
total NNZ (the factor that determines the total number of
\texttt{NZ\_FETCH} and \texttt{MAC} cycles) is held fixed across the
three patterns. The cycle count is therefore dominated by total NNZ
rather than by the spatial distribution of NNZ across rows. At much
larger scales, load imbalance could become visible (very dense rows
would force many sequential $B$-segment loads before the next row begins),
but this effect does not appear at the current benchmark size range. The
Cortex-M0 numbers, by contrast, show small but nonzero differences
between the three patterns (104,488 vs 99,748 vs 102,209 cycles),
indicating that branch prediction and instruction-cache behaviour on
that processor respond slightly to pattern structure in a way the
hardware accelerator does not.

\section{Resource and Timing Evaluation}
\label{sec:resources}

Table~\ref{tab:quartus-final} gives the estimated resource utilisation
for the three build variants. Final Quartus synthesis reports were still
being finalised at the time of writing; the figures below reflect
pre-place-and-route estimates from the Quartus Fitter combined with
known design parameters, and are expected to be replaced by exact
post-fit numbers in the final revision. Utilisation is modest at roughly
16\% of available ALMs, 5\% of M10K BRAMs, and 9\% of DSP blocks, and
leaves substantial room for future expansion. For example, widening the
tile to $TN = 16$ would double the DSP and C-tile register usage and
still remain within the device's capacity. The M10K BRAM count is fixed
at 16 blocks across all three stages because it is determined entirely
by the size of the shared 64\,KB RAM, not by the accelerator's internal
buffers, which sit in fabric registers; and the DSP count of roughly 8
simply reflects the 8 parallel MAC lanes of the $TN = 8$ datapath.

\begin{table}[htbp]
  \centering
  \small
  \begin{tabular}{lccccc}
  \toprule
  Build & ALMs & Registers & M10K BRAMs & DSP blocks & Fmax (MHz) \\
  \midrule
  Baseline accelerator  & $\approx$4,800 & $\approx$7,500 & 16 & 8 & 50+ \\
  Parallel DSP MAC      & $\approx$5,000 & $\approx$7,800 & 16 & 8 & 50+ \\
  INT8 packed (current) & $\approx$5,200 & $\approx$8,000 & 16 & 8 & 50+ \\
  \midrule
  Available (Cyclone~V) & 32,070 & 128,280 & 397 & 87 & --- \\
  \bottomrule
  \end{tabular}
  \caption{Estimated resource utilisation for the three build variants on
           the Cyclone~V 5CSEMA5F31C6. All three designs meet the 50\,MHz
           timing constraint. Final post-fit numbers to be updated in the
           next revision.}
  \label{tab:quartus-final}
\end{table}

\section{Critical Evaluation}

The three-stage experiment provides a clear, quantitative demonstration
that memory bandwidth is the dominant performance bottleneck in CSR SpMM
on this platform, not arithmetic throughput. The $1.32\times$ improvement
from DSP MAC (an arithmetic improvement), set against the $2.18\times$
improvement from INT8 packing (a memory-traffic improvement), is direct
experimental evidence of this conclusion. It is also consistent with the
theoretical analysis in Section~2.2: the operational intensity of CSR
SpMM with $TN = 8$ 32-bit operands is low, placing the design below the
roofline's compute--bandwidth crossover point, and INT8 packing shifts
the operational intensity upward by reducing operand fetch traffic,
moving the design closer to the compute-bound region.

Comparing the $56.80\times$ peak speedup directly against reported
results from OuterSPACE~\citep{outerspace2018}, SpArch~\citep{sparch2020},
or SIGMA~\citep{sigma2020} would not be meaningful: those systems target
orders-of-magnitude larger SpMM workloads on dedicated memory hierarchies.
The relevant question is whether a tightly integrated, MMIO-driven FPGA
accelerator can deliver useful speedup on a constrained embedded
platform, and the answer, based on the measured $56.80\times$ peak and
$45.28\times$ average over PicoRV32 and up to $11.92\times$ over
Cortex-M0, is clearly affirmative. For a wider frame of reference, Bell
and Garland~\citep{bell_garland2009} report that GPU-based CSR SpMV
achieves $2$--$10\times$ speedup over scalar CPU code depending on
matrix structure; the present accelerator, with far lower hardware
complexity and on a constrained SoC, records higher speedups because its
software baseline is a slow in-order core and because INT8 packing
dramatically reduces the memory bottleneck.

Three caveats should be carried forward alongside these results. The
PicoRV32 software baseline is a simple in-order core with no cache, and
against a modern superscalar CPU with vector SIMD extensions the
speedup would be considerably smaller. All data resides on-chip, and a
design with off-chip DRAM would show different numbers because DRAM
latency would affect CPU and accelerator in different proportions.
Finally, the benchmark matrices are synthetic, and real-world GNN or
recommender-system matrices may exhibit structure (extremely uneven row
densities, very high dimensions) that would stress the current
architecture differently. Within these caveats, the results are
internally consistent, reproducible, and supported by a credible
methodology, and the correctness verification at every step confirms
that the recorded speedup numbers reflect real computation rather than
artefacts of incorrect output.
